% Transcription — fonds Alexandre Grothendieck, shelfmark 134-1, batch 1
% (pages 1-6). Established from the facsimile
% archives/batches/134-1/batch-01.pdf (a mirror of
% https://grothendieck.umontpellier.fr/134-1.pdf), per the skill
% transcribe-grothendieck.
%
% Six pages. Page 1 is the folder's own pink cover, carrying only the title
% "Lettres Breen 75, 76" from which the inventory takes its own; page 2 is
% blank. Neither is transcribed, and their numbers do not appear. Pages 3 to 6
% are a single letter, complete, in Grothendieck's hand.
%
% \page{} runs 1 to 6, the position in the folder part, which is what the
% inventory counts and what the site's facsimile pane indexes. The sheets
% themselves carry 60 to 64 in the archivists' pencil — box-level numbering
% for folder 134 taken whole, of which 134-1 is the first part. A note at page
% 3 records this, since a citation of the physical sheet will use the other
% number.
%
% The letter is headed "Lodève le 27.8.74". It is addressed
% "Cher Deligne". Grothendieck says he may run a small algebra seminar on
% "les fourbis de Mme Sinh" — Hoàng Xuân Sính's thesis on Gr-categories, whose
% typescript is folder 135 — transposed into the language of champs, and that
% this led him to the following, which he calls heuristic for the moment.
%
% For M, N abelian sheaves on a topos X he sets E(M,N) = tau_{<=2} RHom(M,N)
% and posits a distinguished triangle (T) relating it to a complex E'(M,N) and
% to Hom(M, 2N)[-2], where 2N is the 2-torsion; page 4 draws from it the exact
% sequence (*), 0 -> Ext^2(M,N) -> H^2(E'(M,N)) -> Hom(M, 2N) -> 0, whose
% middle term he names P(M,N). The heuristic reading is that E'(M,N) expresses
% the strict Picard 2-champ built from Picard 1-champs not necessarily strict,
% and E(M,N) the strict ones. Page 5 states what he can prove — every
% homomorphism M -> 2N comes from a suitable Picard champ, the a priori
% obstruction in Ext^3 being killed by a universal argument. Page 6 applies it
% to K^0 and K^1 of the additive champ of projective A-modules, and closes with
% the question the letter exists for: are (T) and (*) known to the competent —
% he names Quillen, Breen, Illusie — and are there higher variants?
%
% One thing on page 3 is not Grothendieck's. Running vertically up the left
% margin, in another hand, is a covering note: "Cher Breen, j'imagine que tu
% sais ce que le maître demande, mais n'ai pu néanmoins lui répondre. Qu'en
% penses-tu. Bien à toi, P. Deligne." This folder is therefore Breen's copy of
% the letter, forwarded by Deligne, which is what the folder's title records.
% It is transcribed as a marginal note, with a note saying whose it is.
%
% The document is a photocopy and the hand is small and fast; the fine strokes
% are gone in places and a good deal of the connective prose is left as \ill{}
% rather than guessed. The mathematical statements are legible.
%
% The dating is the inventory's own, for the folder.
%
% Folder 103 (pages 8-15) holds another reproduction of the same leaves, the
% same strokes line for line, sharper on the dark strokes and four pages
% longer. Its fine strokes are paler than here, so the two copies were read
% against each other; three readings of this file were revised on 2026-09-13
% by Opus 5 (claude-opus-5), and each only where this folder's own facsimile
% bears the reading out. (1) The date: the first pass read "27." followed by a
% struck-over month and year. The bar over "8.74" is the long top stroke of
% the 7 of "74", thickened by the photocopy; the digits are legible here, and
% the 103 copy shows them with no deletion. It is now "27.8.74", with "8.74"
% marked uncertain. (On the 103 copy the "2" has all but faded, and that file
% reads "7.8.74"; it is not edited from here.) (2) Page 5: "s'obtient en" is
% not struck but written between the lines, with a bracket before
% "associant" marking where it goes; the stroke under it runs below the
% letters, not through them. The sentence reads "l'invariant qui s'obtient en
% associant à toute section L", L also being an interlinear insertion. (3) The
% label of the exact sequence is "(*)", not "(S)": it is faint at page 4's
% display, but legible in the text of page 4 ("La suite exacte (*)") and of
% page 6 ("de (*)"). A fourth followed on the same day: page 5 had
% "\struck{splittée}", but on both copies "splittée :" stands and the
% deletion is the short word after the colon, now \struck{\ill{}}.
%
% Pass: Opus 5 (claude-opus-5), 2026-08-23 — first pass, unchecked against the pages by a human.
\documentclass[11pt,a4paper]{article}
% Le préambule commun à toutes les transcriptions du fonds Grothendieck.
%
% Il tient en une idée : **l'incertitude se compose, elle ne se résout pas.**
% Une transcription qui lisse un mot douteux a détruit la seule chose pour
% laquelle on la faisait — un lecteur doit pouvoir savoir, à chaque endroit,
% ce qui était sur le feuillet et ce qui a été ajouté. Les six macros qui
% suivent sont donc l'appareil critique, et non de la décoration.
%
% Les mêmes commandes sont reconnues par `scripts/render.mjs`, qui produit la
% vue de lecture du site. Une macro ajoutée ici doit l'être aussi là-bas,
% faute de quoi le rendu échouera bruyamment — ce qui est voulu : un rendu
% silencieusement faux est pire qu'un rendu absent.

\ProvidesPackage{grothendieck}[2026/08/08 Transcriptions du fonds Grothendieck]

\usepackage{amsmath,amssymb,amsthm}
\usepackage{mathrsfs}
\usepackage{stmaryrd}
% Les diagrammes commutatifs sont partout dans ce fonds : c'est la langue
% même de la géométrie algébrique de Grothendieck, et une transcription qui
% les rendrait en prose ne transcrirait rien.
\usepackage{tikz-cd}
% Français par défaut — c'est la langue des feuillets — mais l'anglais est
% chargé aussi : la traduction et le résumé sont des documents anglais, et la
% typographie française (espaces avant les deux-points) y serait une faute.
% Ces documents basculent par \selectlanguage{english} en tête de corps.
\usepackage[english,french]{babel}
\usepackage{fontspec}
\usepackage{geometry}
\geometry{a4paper,margin=25mm,marginparwidth=28mm}
\usepackage{marginnote}
\usepackage{xcolor}
% ulem, not soul. `soul`'s \st and \ul are built on a character-by-character
% scan that XeLaTeX's font machinery breaks: the document fails with an
% undefined control sequence rather than degrading. ulem does the same two jobs
% and is Unicode-engine safe. `normalem` stops it hijacking \emph, which the
% transcriptions use for Grothendieck's own emphasis.
\usepackage[normalem]{ulem}
\usepackage{hyperref}
\hypersetup{colorlinks=true,linkcolor=black,urlcolor=black,pdfborder={0 0 0}}

% --- Métadonnées du lot -----------------------------------------------------
% Reprises telles quelles par le script de rendu, qui les affiche en tête de
% la vue de lecture. Elles fixent ce que le fichier prétend couvrir : sans
% elles, une transcription flotte sans point d'attache dans un fonds de seize
% mille feuillets.

\newcommand{\folder}[1]{\gdef\grothFolder{#1}}
\newcommand{\batch}[1]{\gdef\grothBatch{#1}}
\newcommand{\leaves}[2]{\gdef\grothFirst{#1}\gdef\grothLast{#2}}
\newcommand{\foldertitle}[1]{\gdef\grothFoldertitle{#1}}
\gdef\grothFolder{?}\gdef\grothBatch{?}\gdef\grothFirst{?}\gdef\grothLast{?}\gdef\grothFoldertitle{}

% --- Le résumé --------------------------------------------------------------
%
% Il ouvre la lecture modernisée, sous le titre « Résumé », et il est composé
% en petit. Ce n'est pas un ornement : c'est le seul passage du document qui
% ne soit pas des mathématiques, et un lecteur doit pouvoir le reconnaître
% avant de l'avoir lu — puis le sauter s'il connaît déjà le sujet, ou s'y
% arrêter s'il ne le connaît pas. Le filet qui le suit marque la reprise du
% corps.

\newenvironment{resume}
  {\small\setlength{\parskip}{0.45em}}
  {\par\medskip\hrule\medskip}

% --- Mots-clés ---------------------------------------------------------------
%
% En anglais, en clôture du résumé de la lecture modernisée : le vocabulaire
% moderne sous lequel chercher ce que les feuillets construisent. Le manifeste
% du site (`npm run manifest`) les extrait pour étiqueter la cote entière —
% cette ligne est la source unique des tags, il n'y a pas de fichier à côté.
\newcommand{\keywords}[1]{\par\smallskip\noindent
  {\footnotesize\textsc{Keywords} --- \textit{#1}}\par}

% --- L'appareil critique ----------------------------------------------------

% Le numéro de feuillet, en marge. C'est l'ancre qui permet de régler un
% désaccord sur une formule en regardant *un* feuillet plutôt qu'une chemise
% de six cents. Le site s'en sert aussi pour tourner les pages du fac-similé
% quand on fait défiler la transcription.
\newcommand{\leaf}[1]{%
  \marginnote{\footnotesize\color{gray!70}\textsf{#1}}%
  \label{leaf:#1}\ignorespaces}

% Illisible : on ne devine pas. Un mot inventé qui se lit comme les autres est
% le pire résultat possible.
\newcommand{\ill}{\textcolor{red!60!black}{[\,\ldots\,]}}

% Lecture douteuse : le mot est donné, et signalé comme incertain.
\newcommand{\uncertain}[1]{\uline{#1}}

% Ajout de l'éditeur — une lettre manquante, un mot sous-entendu.
\newcommand{\add}[1]{\textcolor{blue!50!black}{[#1]}}

% Biffé par Grothendieck lui-même. Ce qu'il a rayé fait partie du document :
% une rature dit souvent où la pensée a changé de direction.
\newcommand{\struck}[1]{\sout{#1}}

% Note du transcripteur, sur l'état matériel du feuillet.
\newcommand{\note}[1]{\par\smallskip{\small\color{gray!60!black}\textit{[#1]}}\par\smallskip}

% Note marginale de Grothendieck, distincte du corps du texte.
\newcommand{\marginal}[1]{\par\smallskip\noindent\hspace{1em}%
  {\small\color{gray!60!black}#1}\par\smallskip}

% --- Titre ------------------------------------------------------------------

\AtBeginDocument{%
  \begin{center}
    {\large\bfseries Fonds Alexandre Grothendieck --- cote n\textsuperscript{o}~\grothFolder}\\[2pt]
    {\itshape\grothFoldertitle}\\[4pt]
    {\small feuillets \grothFirst--\grothLast\ (lot \grothBatch)}\\[2pt]
    {\footnotesize\color{gray!60!black}Transcription établie d'après le fac-similé numérisé
     par l'Université de Montpellier.\\
     Les lectures incertaines sont soulignées, les illisibles notées [\,\ldots\,],
     les ajouts entre crochets.}
  \end{center}
  \vspace{1em}}

\endinput


\folder{134-1}
\batch{1}
\pages{1}{6}
\dating{1974-1984}
\watermark{Édition de démonstration}
\foldertitle{Lettres Breen 75,76 [correspondance avec Larry Breen, Pierre Deligne et Daniel Quillen] : lettres et copies de lettre (1974-1976, 1984)}

\begin{document}

\section*{A.\ Grothendieck à P.\ Deligne, de Lodève (pages 3 à 6)}

\page{3}
\note{Les feuillets portent 60 à 64 au crayon de l'archiviste : c'est la
numérotation de la cote 134 prise entière, dont 134-1 est la première partie.
Les numéros employés ici, 1 à 6, sont ceux que compte l'inventaire pour cette
partie.}

\marginal{Cher Breen, j'imagine que tu sais ce que le maître demande, mais
n'ai pu néanmoins lui répondre. Qu'en penses-tu. Bien à toi, P.\ Deligne}

\note{Cette note court verticalement dans la marge de gauche, d'une autre
main que celle de la lettre, et elle est signée. Le dossier est donc
l'exemplaire de Breen, que Deligne lui a fait suivre.}

Lodève le 27.\uncertain{8.74}\note{Le trait qui court sur « 8.74 » est la
barre du 7 de « 74 », épaissie par la photocopie, et non une rature :
l'exemplaire de cette lettre que conserve le dossier 103 (page 8) montre ces
chiffres nets et sans surcharge.}

Cher Deligne,

Étant peut-être empêché \ill{} \ill{} faute d'\uncertain{assurer} un cours de
1er cycle au 1er trimestre, je vais
peut-être \ill{} faire un petit séminaire d'algèbre, et envisage de le faire
sur les fourbis de Mme Sinh, éventuellement transposés dans le contexte des
« champs ». À ce propos, je tombe sur le truc suivant, qui pour l'instant
\uncertain{reste} heuristique. Si $M$, $N$ sont deux faisceaux abéliens sur
un topos $X$, et si

\[ \tau_{\leq 2} \, R\underline{\mathrm{Hom}}(M, N) = E(M, N) \]

est le complexe ayant les invariants

\[ \begin{cases}
\underline{H}^i = \underline{\mathrm{Ext}}^i(M, N) & \text{pour } 0 \leq i \leq 2 \\
\underline{H}^i = 0 & \text{si } i \notin [0, 2]
\end{cases} \]

il doit y avoir un triangle distingué canonique

\begin{tikzcd}
& \underline{\mathrm{Hom}}(M, {}_2 N)[-2] \arrow[dl] & \\
E(M, N) \arrow[rr] & & E'(M, N) \arrow[ul]
\end{tikzcd}

\[ (T) \]

\page{4}
donc $E'(M,N)$ est un complexe dont les invariants $\underline{H}^i$ sont ceux
de $E(M,N)$ en degré $i+2$, et qui en degré 2 donne lieu à une suite exacte

\[ (*) \qquad 0 \longrightarrow \underline{\mathrm{Ext}}^2(M, N)
   \longrightarrow \underbrace{\underline{H}^2(E'(M,N))}_{P(M,N)}
   \xrightarrow{\ \sigma\ } \underline{\mathrm{Hom}}(M, {}_2 N)
   \longrightarrow 0 \]

Heuristiquement, $E'(M,N)$ est le complexe qui exprime le « 2-champ de Picard
"strict" » formé des 1-champs de Picard (pas nécessairement stricts)
« épinglés » par $M$, $N$, sur des objets variables de $X$, en admettant que
la théorie pour les 1-champs de Picard stricts s'étend aux 2-champs de Picard
stricts (ce qui pour moi ne fait guère de doute) ; de même $E(M,N)$ correspond
aux champs de Picard stricts épinglés par $M$, $N$.

La suite exacte $(*)$ se construit en tout cas canoniquement « à la main », où
le terme médian est le faisceau des \uncertain{données} à « équivalence »
près des

\page{5}
\struck{catégories} champs de Picard épinglés par $M$, $N$, $\sigma$ étant
l'invariant qui s'obtient en associant à toute section
$L$\note{« s'obtient en » et « $L$ » sont écrits en interligne, un crochet
devant « associant » marquant la place du premier ; le trait sous « s'obtient
en » passe sous les lettres et ne les biffe pas.} d'un champ de Picard, la symétrie de \struck{\ill{}} $L \otimes L$,
interprétée comme section de ${}_2 N$. Je sais prouver (sauf erreur) que tout
hom.\ $M \to {}_2 N$ provient d'un champ de Picard convenable (épinglé par
$M$, $N$) (a priori l'obstruction est dans $\mathrm{Ext}^3(X; M, N)$, mais un
argument « universel » prouve qu'elle est nulle). Cela prouve que l'extension
$(*)$ est bien proche d'être splittée : \struck{\ill{}} toute section du troisième
faisceau, sur un objet quelconque de $X$, se relève --- en d'autres
\struck{\ill{}} termes l'extension a une section « ensembliste ». Bien sûr,
il y a mieux en fait : toute section sur un $U \in \mathrm{Ob}\, X$
« provient » d'un élément de $\mathbb{H}^2(U, E'(M,N))$ (hypercohomologie,
\uncertain{i.e.} $\mathbb{H}^2$).

\page{6}
\ill{} soit $A$ un anneau sur $X$, soient $M$, $N$ respectivement les
faisceaux $K^0$ et $K^1$ associés au champ additif des $A$-Modules projectifs
\ill{} (p.\ ex.). Alors le raisonnement de Mme Sinh nous fournit un champ de
Picard épinglé par $M$, $N$, d'où une section canonique du terme médian
$P(M,N)$ de $(*)$.

\emph{NB} Tout ce qui précède a les fonctorialités évidentes en $M$, $N$,
$X$ …

\emph{Question} : Le triangle exact \uncertain{$(T)$} et la suite exacte
$(*)$ sont-ils connus par les compétents (Quillen, Breen, Illusie…) ?
Connaissent-ils des variantes « supérieures » ? (Un principe
« géométrique » pour les obtenir pourrait être via des $n$-champs de Picard
non néc.\ stricts…)

\end{document}
